\chapter{Regrouping isomorphism for the affine base-change tensor tower}
\label{chap:Cohomology_RegroupHelper}
% This chapter blueprints the standalone, geometry-free helper that lives in
% AlgebraicJacobian/Cohomology/RegroupHelper.lean (1:1 slug). It is the pure
% tensor-algebra core of lem:base_change_mate_regroupEquiv in the FlatBaseChange
% chapter; it is split into its own compilation unit so the object-level
% ModuleCat iso can be obtained by a single LinearEquiv.toModuleIso application
% (the separate import normalises the `Module A (A ⊗[R] R')` instance diamond
% that blocks the same-file construction).

\section{The R'-linear regrouping equivalence}

\begin{lemma}

\leanok
[Regrouping linear equivalence for the tensor tower]
  \label{lem:base_change_regroup_linearEquiv}
  \lean{AlgebraicGeometry.base_change_regroup_linearEquiv}
  \uses{lem:isPushout_cancelBaseChange_mathlib}
  % SOURCE: [Stacks Project], Cohomology of Schemes, Lemma
  % `lemma-affine-base-change' (\S\,Cohomology and base change, I), proof, the
  % "boils down to the equality" step (read from references/stacks-coherent.tex,
  % L933--938).
  % SOURCE QUOTE: "Thus we see that the lemma boils down to the
  % equality
  % $$
  % (R' \otimes_R A) \otimes_A M = R' \otimes_R M
  % $$
  % as $R'$-modules."
  \textit{Source: Stacks Project, Cohomology of Schemes, Lemma ``Affine base
  change'' (proof, ``boils down to the equality'' step).}
  Let \(\psi : R \to R'\) and \(\varphi : R \to A\) be ring homomorphisms and \(M\)
  an \(A\)-module. On the \emph{canonical} Mathlib instances (with \(R'\) acting on
  \(A \otimes_R R'\) through the \(\texttt{Algebra.TensorProduct.rightAlgebra}\)
  factor, i.e.\ \(r' \cdot (a \otimes s) = a \otimes (r' s)\)) there is a bundled
  \(R'\)-linear isomorphism
  \[
    \operatorname{regroup} :
    (A \otimes_R R') \otimes_A M \;\xrightarrow{\ \sim\ }\; R' \otimes_R M ,
  \]
  the ``equality as \(R'\)-modules'' of the affine base-change source. On generators
  it sends \((a \otimes r') \otimes m \mapsto r' \otimes (a \cdot m)\); in
  particular \((1 \otimes r') \otimes m \mapsto r' \otimes m\), with inverse
  \(r' \otimes m \mapsto (1 \otimes r') \otimes m\). It carries \emph{no} flatness
  hypothesis.
\end{lemma}

\begin{proof}
  \uses{lem:isPushout_cancelBaseChange_mathlib}
  \leanok
  This is pure tensor algebra, and it is supplied \emph{directly} by Mathlib's
  pushout-form base-change cancellation
  \(\operatorname{Algebra.IsPushout.cancelBaseChange}\)
  (Lemma~\ref{lem:isPushout_cancelBaseChange_mathlib}), instantiated at \(S = R'\),
  \(B = A \otimes_R R'\):
  \[
    \operatorname{cancelBaseChange} R\,R'\,A\,(A \otimes_R R')\,M :
    (A \otimes_R R') \otimes_A M \;\xrightarrow{\ \sim\ }\; R' \otimes_R M .
  \]
  The required pushout instance \(\operatorname{Algebra.IsPushout} R\,R'\,A\,
  (A \otimes_R R')\) is supplied for free by \(\operatorname{TensorProduct.isPushout'}\).
  The decisive structural point is that this equivalence is \emph{natively}
  \(R'\)-linear: the cancellation is bundled as a \(\operatorname{LinearEquiv}\) over
  \(R'\) (its base-ring slot is \(S = R'\)), so the \(R'\)-linearity is part of the
  Mathlib datum with no hand-written scalar-compatibility obligation.
  On the generator it computes, by \(\texttt{cancelBaseChange\_tmul}\),
  \[
    (a \otimes r') \otimes m \;\longmapsto\; r' \otimes (a \cdot m) ,
  \]
  with inverse \(r' \otimes m \mapsto (1 \otimes r') \otimes m\)
  (\(\texttt{cancelBaseChange\_symm\_tmul}\)); the \(R'\)-action enters through the
  \(A \otimes_R R'\) factor (\(r' \cdot (a \otimes s) = a \otimes (r' s)\)). Being a
  single bundled Mathlib linear equivalence it is an \(R'\)-linear isomorphism, and it
  carries no flatness hypothesis.

\end{proof}
